\chapter{Guía rápida de comandos forenses en consola GNU/Linux}

Este apéndice constituye una guía de referencia rápida (\textit{cheat sheet}) con los comandos esenciales que todo analista e investigador forense debe dominar al interactuar con sistemas basados en Unix y GNU/Linux, tanto en escenarios de respuesta en vivo (\textit{live response}) como en el examen estático de imágenes (\textit{dead-box forensics}).

\section{Inspección y triage del sistema en vivo}

\subsection{Procesos en ejecución y memoria volátil}
\begin{lstlisting}[style=terminal]
# Listar todos los procesos con usuario, PID, consumo y linea de comando completa:
ps auxwww

# Arbol jerarquico de procesos padre-hijo (revela procesos huerfanos o inyectados):
pstree -p -a

# Descriptores de archivos abiertos y conexiones por proceso (ej. PID 1234):
lsof -p 1234

# Verificacion del binario real en disco que ejecuta un proceso:
ls -l /proc/1234/exe

# Inspeccion de variables de entorno de un proceso sospechoso:
strings -a /proc/1234/environ | tr '\0' '\n'

# Volcado del mapa de memoria virtual de un proceso:
cat /proc/1234/maps
\end{lstlisting}

\subsection{Conexiones de red y sockets activos}
\begin{lstlisting}[style=terminal]
# Sockets TCP y UDP en escucha y establecidos con PID y nombre de proceso:
netstat -tulpn
# Alternativa moderna (iproute2):
ss -tulpan

# Listar procesos que mantienen conexiones de red activas:
lsof -i -n -P

# Captura de trafico crudo en una interfaz especifica (sin resolucion DNS):
tcpdump -i eth0 -nn -s 0 -w captura_incidente.pcap
\end{lstlisting}

\subsection{Búsqueda de binarios privilegiados y atributos especiales}
\begin{lstlisting}[style=terminal]
# Localizar todos los ejecutables con bit SUID o SGID en el sistema:
find / -type f \( -perm -04000 -o -perm -02000 \) -exec ls -l {} \; 2>/dev/null

# Buscar archivos pertenecientes a usuarios no estandar en directorios de sistema:
find /bin /sbin /usr/bin /usr/sbin /dev -not -user root -ls

# Auditoria recursiva de atributos inmutables (+i) y append-only (+a) en Ext2/3/4:
lsattr -aR /dev /bin /sbin /usr/lib /etc 2>/dev/null | grep -E "(\-i\-|\-a\-)"
\end{lstlisting}

\section{Análisis de imágenes con The Sleuth Kit (TSK)}

Cuando se analiza una copia binaria de un disco o partición (\texttt{disco.img}):

\begin{lstlisting}[style=terminal]
# 1. Informacion general del sistema de archivos (tamaño de bloque, inodos):
fsstat disco.img

# 2. Listado recursivo de archivos con inodos asociados (incluyendo borrados):
fls -r -p -m "/" disco.img > bodyfile.txt

# 3. Inspeccion detallada de un inodo especifico (marcas MAC, bloques asignados):
istat disco.img 533487

# 4. Extraccion directa del contenido de un inodo hacia un archivo local:
icat disco.img 533487 > binario_rescatado.bin

# 5. Listar unicamente inodos eliminados / huerfanos:
ils -r -m disco.img >> bodyfile.txt

# 6. Generacion de la linea de tiempo cronologica formateada (mactime):
mactime -b bodyfile.txt -d -z CST > linea_tiempo.csv
\end{lstlisting}

\section{Verificación de integridad y detección de alteraciones}

\begin{lstlisting}[style=terminal]
# Calculo de huellas hash criptograficas para preservacion de custodia:
md5sum disco.img > disco.img.md5
sha256sum disco.img > disco.img.sha256

# Extraccion de cadenas de texto ASCII y Unicode de un archivo sospechoso:
strings -a -t x archivo_sospechoso.bin

# Auditoria de integridad de paquetes en distribuciones Debian/Ubuntu:
debsums -c -s

# Auditoria de integridad de paquetes en distribuciones Red Hat/CentOS/Fedora:
rpm -Va
\end{lstlisting}
